\section{CONCLUSIONES}

El desarrollo del buscador semántico permitió comprobar que las tecnologías de la Web
Semántica resultan adecuadas para modelar dominios de alojamiento temporal donde la
información relevante no puede reducirse a filtros aislados. La ontología construida en
OWL/RDF representó de forma coherente propiedades, amenidades, perfiles de huésped,
propósitos de viaje, políticas y contexto geográfico, habilitando consultas más
expresivas que las disponibles en un catálogo convencional.

La integración con DBpedia cumplió un doble propósito. Por un lado, aportó una base
externa confiable para enlazar tipos de alojamiento, amenidades y ciudades bolivianas
mediante \texttt{owl:sameAs} y \texttt{owl:equivalentClass}. Por otro lado, permitió
justificar la interoperabilidad del proyecto con el ecosistema de \textit{Linked Data}
sin comprometer la estabilidad del prototipo, dado que la operación principal se ejecuta
sobre un dataset local en Fuseki.

La implementación de multilingüidad mediante \texttt{rdfs:label} y
\texttt{rdfs:comment} demostró ser una decisión técnicamente sólida. A partir de un
único grafo RDF fue posible publicar y recuperar resultados en español, inglés y
francés, evitando duplicación de entidades y manteniendo una representación semántica
consistente entre idiomas. Esta estrategia confirma la ventaja de RDF/OWL para soportar
internacionalización de forma nativa.

Desde la perspectiva de ingeniería, el uso de Fuseki como \textit{triple store}, SPARQL
como lenguaje de consulta y Next.js como capa de presentación permitió construir un
prototipo funcional y cercano a un escenario real de uso. El sistema no quedó limitado
a la teoría ontológica: ofrece interfaz web, resultados visuales, navegación por idioma
y un flujo de búsqueda operable por el usuario final.

La validación con HermiT y las pruebas funcionales realizadas muestran que la solución es
coherente tanto en el plano lógico como en el operativo. En consecuencia, se concluye
que el objetivo de la práctica fue alcanzado: se desarrolló un buscador semántico basado
en OWL/RDF, enriquecido con DBpedia, con soporte multilingüe y materializado en una
aplicación de software defendible académicamente.

Como trabajo futuro, el proyecto puede ampliarse con nuevas ciudades, más puntos de
interés, consultas geográficas basadas en coordenadas, reglas de razonamiento más ricas
y un módulo de administración de datos que permita poblar la ontología de forma más
automatizada. Sin embargo, aun en su estado actual, el prototipo constituye una prueba
completa de aplicación de técnicas de Web Semántica en un problema concreto del dominio
de alojamiento temporal.
